%% ============================================================
%% 片段：s03a-quotient   对应原文 PDF 第 32--47 页
%% 本文件由 gw/convert.py 从 arXiv 源码（hep-th/0612073v2）自动生成：
%%   * 公式/编号/引用/图 已经锁定为原文内容，请勿改动任何数学内容与编号；
%%   * 只需把英文正文译为中文，并在每节开头加“本节重点”、关键处加译注。
%% ============================================================


\section{经典几何续论}

%% ==== 原文 §3：PDF 第 32 页（印刷第 31 页）====
\label{sec:3}

\begin{keytakeaways}
  \begin{itemize}[leftmargin=1.4em,itemsep=2pt,topsep=2pt]
    \item \textbf{物理图像}：带分歧的希格斯丛模空间 $\MH(\alpha,\beta,\gamma;p)$ 是超 Kähler 商 $\EUW(\alpha,\beta,\gamma;p)/\neg/\neg/\EUG_p$；参数 $(\alpha,\beta,\gamma)\in\mathfrak t\times\mathfrak t\times\mathfrak t$ 刻画面算符在分歧点 $p$ 处的奇异类型，而面算符所对应的量子场论问题由列维子群 $\mathbb{L}$ 的选择决定。
    \item \textbf{数学工具}：超 Kähler 矩映射 $\vec\mu=(\mu_I,\mu_J,\mu_K)$ 与三个 Kähler 形式 $\omega_I,\omega_J,\omega_K$；复观点在适当的稳定性条件下把超 Kähler 商等同于几何不变量理论商 $\nu_I^{-1}(0)/\EUG_{p,\C}$，从而与稳定抛物希格斯丛（复结构 $I$）或稳定抛物局部系统（复结构 $J$）的模空间联系起来。
    \item \textbf{关键结论}：三个 Kähler 形式的上同调类都是 $(\alpha,\beta,\gamma)$ 的线性函数，即式~\eqref{eq:3.8}--\eqref{eq:3.10}；并且只要避开实余维至少为二的奇点，实辛簇 $(\MH,\omega_K)$ 就与 $\alpha$、$\beta$ 无关。
    \item \textbf{非正则情形}：$(\alpha,\gamma)$ 或 $(\beta,\gamma)$ 非正则时，单值性只能确定到共轭类的闭包中：半单部分由 $U=\exp(-2\pi(\alpha-i\gamma))$ 给出，幂幺部分落在由 $\beta$ 决定的幂幺根 $\EUN$ 内，一般情形下落在理查森类中；相应的局部解由纳姆方程（Nahm's equations）描述（脚注~17、第~\ref{sec:3.8}~节）。
  \end{itemize}
\end{keytakeaways}

本节将重新考察第~\ref{sec:2.2}~节的构造，并描述那些需要更深入地理解 $\MH$ 的经典几何才能说明的方面。我们力求给出一个相对简单——但当然并不数学上完整——的介绍，涵盖那些目前或可能对几何朗兰兹纲领有用的内容。我们所讨论的大部分课题在数学文献中已有远为充分的处理，但有几处也许是新的，特别是仿射外尔群（affine Weyl group）在 $\MH$ 的上同调（cohomology）上的作用（第~\ref{sec:3.7}~节）、辛形式的上同调类的线性性（第~\ref{sec:3.2}~节），以及 $\MH$ 在非正则点处局部奇点的性质（第~\ref{sec:3.6}~节）。

我们将描述的许多最重要现象都有局部模型，这些模型涉及复余伴随轨道（complex coadjoint orbit）上的超 Kähler（hyper-Kähler）度规，其构造见 \cite{ref34,ref35,ref36,ref37}。我们把这部分推迟到第~\ref{sec:3.8}~节。



\subsection{超 Kähler 商}
 %% ==== 原文 §3.1：PDF 第 32 页（印刷第 31 页）====
\label{sec:3.1}

黎曼面 $C$ 上希格斯丛（Higgs bundle）的模空间（moduli space）$\MH(G,C)$ 的最基本性质之一是，它可以构造为超 Kähler 商（hyper-Kähler quotient）\cite{ref21}。我们首先回顾无奇点时的构造。取光滑的 $G$-丛 $E\to C$，以 $\EUW$ 记数对 $(A,\phi)$ 构成的空间，其中 $A$ 是 $E$ 上的联络（connection），$\phi\in
\Omega^1(C,{\rm ad}(E))$。$\EUW$ 是平坦的超 Kähler 流形；详细论述见 \cite{ref8} 第 4.1 节。设 $\EUG$ 是 $E$ 的规范变换（gauge transformation）群。则 $\EUG$ 作用在 $\EUW$ 上，保持超 Kähler 结构，并带有超 Kähler 矩映射（moment map）$\vec\mu=(\mu_I,\mu_J,\mu_K)$：

% 式 (3.1)
\setcounter{equation}{0}
\begin{equation}\label{eq:3.1}
  \begin{aligned}\mu_I & =-{1\over
2\pi}\int_C\,\Tr\,\epsilon(F-\phi\wedge \phi),\\ \mu_J &
=-{1\over 2\pi}\int_C|d^2z|\,\Tr\,\epsilon\left(D_z\phi_{\bar
z}+D_{\bar z}\phi_z\right), \\ \mu_K& = -{i\over
2\pi}\int_C|d^2z|\,\Tr\,\epsilon\left(D_z\phi_{\bar z}-D_{\bar
z}\phi_z\right).\end{aligned}
\end{equation}
这里 $\epsilon\in \Omega^0(C,{\rm ad}(E))$
是 $\EUG$ 的李代数的一个元素；$\vec\mu$ 关于 $\epsilon$ 是线性的。给定超 Kähler 流形 $\EUW$ 以及保持超 Kähler 结构、作用在 $\EUW$ 上并带有矩映射 $\vec\mu$ 的群
$\EUG$，超 Kähler 商
$\EUW/\neg/\neg/\EUG$ 一般地 \cite{ref49} 定义为 $\vec\mu^{-1}(0)/\EUG$。该超 Kähler 商是超 Kähler 流形。在目前的情形下，方程
$\vec\mu=0$（对所有 $\epsilon$）正是希钦方程（Hitchin's equations）\eqref{eq:2.1}。于是超 Kähler 商 $\EUW/\neg/\neg/\EUG$ 恰好是希钦方程之解的模空间 $\MH$。

\begin{yizhu}[易错点]
这里的除法应理解为几何不变量理论意义下的商，而非集合论意义上的商：$\EUG$ 在 $\vec\mu^{-1}(0)$ 上的作用一般有非平凡稳定子，商空间因此会出现轨形（orbifold）奇点。第~\ref{sec:3.2}~节把同一对象改写成复结构 $I$ 下的 GIT 商 $\nu_I^{-1}(0)/\EUG_{p,\C}$ 时，这一点表现得最清楚。GIT 商的稳定性条件一般依赖于 $\alpha$（当 $\sigma$ 一般时与 $\alpha$ 无关，见第~\ref{sec:3.2}~节），所以“超 Kähler 商 $=$ 辛商 $=$ GIT 商”这一串等同并非对任意参数都自动成立；跨墙（wall-crossing）现象正是稳定性条件改变时发生的。另注意 $\MH$ 只通过 $\alpha$ 在 $\mathfrak t/\Lambda_{\cochar}=\TT$ 中的像依赖于 $\alpha$（脚注~14）。
\end{yizhu}

这一构造可以推广到带有熟知类型孤立奇点的希格斯丛 \cite{ref23,ref24}。取一点 $p\in C$，并把丛 $E$ 在 $p $ 处的结构群约化为环面 $\TT$，以 $\EUG_p$ 记 $\EUG$ 中在 $p$ 处的限制落在 $\TT$ 内的规范变换构成的子群。再取三元组\footnote{正如第~\ref{sec:2.1}~节所解释的，$\MH$ 实际上只依赖于 $\alpha$ 在 $\mathfrak t/\Lambda_{\cochar}=\TT$ 中的像。}
$(\alpha,\beta,\gamma)\in \mathfrak t$（$\TT$ 的李代数）。在 $p$ 附近取坐标 $r,\theta$，以 $\EUW(\alpha,\beta,\gamma;p)$ 记在 $p$ 附近具有熟知类型奇点的数对 $(A,\phi)$ 构成的空间：

% 式 (3.2)
\setcounter{equation}{1}
\begin{equation}\label{eq:3.2}
  \begin{aligned}A & =\alpha\,d\theta+\dots\\
 \phi & = \beta\,{dr\over r}-\gamma\,d\theta+\dots.\end{aligned}
\end{equation}

省略号表示比 $1/r$ 更不奇异的项。又记 $\EUG_p$ 为 $E$ 的规范变换中在 $p$ 处取值于 $\TT$ 的群。这些正是保持条件~\eqref{eq:3.2} 不变的规范变换。于是模空间 $\M_H(\alpha,\beta,\gamma;p)$ 就是超 Kähler 商
$\EUW(\alpha,\beta,\gamma;p)/\neg/\neg/\EUG_p$。我们称之为在点 $p$ 处带分歧的希格斯丛的模空间。类似地，可以把在若干点处带分歧的希格斯丛的相应模空间也构造为超 Kähler 商，即在若干点 $p_1,\dots,p_s\in C$ 处具有奇点、并由参数 $(\alpha_i,\beta_i,\gamma_i)\in \mathfrak t$（$i=1,\dots,s$）标记的希钦方程之解的模空间。为记号简洁起见，本节主要就单点分歧的情形陈述，但所有论断对更一般的情形都有直接类比。

具有复结构（complex structure）$I,J,$ 和 $K$ 的超 Kähler 流形，具有相应的 Kähler 形式 $\omega_I$、$\omega_J$ 和
$\omega_K$。在目前的情形下，这些形式可以像无奇点时那样写出，\footnote{与 \cite{ref8} 一样，我们以 $\delta$ 记 $\EUW$ 上的外微分，以 $d$ 记 $C$ 上的外微分。}
% 式 (3.3)
\setcounter{equation}{2}
\begin{equation}\label{eq:3.3}
  \begin{aligned}\omega_I & =
-{i\over 2\pi}\int_C|d^2z|\,\Tr\left(\delta A_{\bar z}\wedge
\delta A_z-\delta \phi_{\bar z}\wedge \delta\phi_z\right)\\
&=-{1\over 4\pi}\int_C\,\Tr\left(\delta A\wedge \delta A
-\delta\phi\wedge\delta \phi\right) \\
                      \omega_J & ={1\over 2\pi}\int_C |d^2z|\,\Tr\left(
                      \delta\phi_{\bar z}\wedge \delta A_z+\delta\phi_z\wedge
                      \delta A_{\bar z} \right)\\
                      \omega_K & =
                      {i\over 2\pi}\int_C|d^2z|\,\Tr\left(\delta\phi_{\bar z}
                      \wedge\delta A_z-\delta\phi_z\wedge\delta A_{\bar z}\right)\\
                      & = {1\over 2\pi}\int_C \,\Tr\,\delta\phi
                      \wedge\delta A.\end{aligned}
\end{equation}


正如无奇点时那样（见 \cite{ref8} 第 4.1 节），$\omega_I$ 与
% 式 (3.4)
\setcounter{equation}{3}
\begin{equation}\label{eq:3.4}
  \omega_I'  =
-{i\over 2\pi}\int_C|d^2z|\,\Tr\left(\delta A_{\bar z}\wedge
\delta A_z\right)=-{1\over 4\pi}\int_C\,\Tr\left(\delta A\wedge
\delta A \right),
\end{equation}
上同调，后者只依赖于 $A$，因而来自 $G$-丛的模空间 $\M$ 的拉回。事实上，令 $\lambda_I={1\over 4\pi}\int_C\Tr
\,\phi\wedge\delta \phi$。$\lambda_I$ 是规范不变的，且（在分部积分并使用希钦方程之后）与规范变换的生成元 $\delta A=-d_A\epsilon,
\,\delta\phi=[\epsilon,\phi]$ 缩并时为零。因此它是 $\MH$ 上某个一形式的拉回，而公式
$\omega_I-\omega_I'=\delta\lambda_I$ 表明 $\omega_I$ 与 $\omega_I'$ 上同调等价。

\begin{yizhu}[易错点]
超 Kähler 流形的三个复结构 $I,J,K$ 与三个 Kähler 形式 $\omega_I,\omega_J,\omega_K$ 一一对应（约定 $\omega_I(\cdot,\cdot)=g(I\cdot,\cdot)$ 等），而三个全纯辛形式是 $\Omega_I=\omega_J+i\omega_K$、$\Omega_J=\omega_K+i\omega_I$、$\Omega_K=\omega_I+i\omega_J$。核对本节结论时须注意指标并不“对齐”：式~\eqref{eq:3.8} 说 $[\,\omega_I]$ 只依赖 $\alpha$，而第~\ref{sec:3.2}~节说 $\Omega_I=\omega_J+i\omega_K$ 的全纯参数是 $\beta+i\gamma$；类似地 $\Omega_J$ 依赖 $\gamma+i\alpha$、$\Omega_K$ 依赖 $\alpha+i\beta$。把这两组对应关系混同是本节最常见的错误来源。
\end{yizhu}

无奇点时，$\omega_J$ 与 $\omega_K$ 的上同调类为零。例如，对 $\omega_K$ 可以写出 $\omega_K=\delta\lambda_K$（其中 $\lambda_K={1\over 2\pi}\int_C\Tr \,\phi\wedge\delta A$）来证明这一点。只要 $\gamma=0$，这在存在奇点时仍然成立；但当 $\gamma\not=0$ 时，$\lambda_K$ 与规范变换的生成元缩并后不为零，因而它不是 $\MH$ 上一形式的拉回。事实上，与规范变换的生成元缩并把 $\omega_K$ 映为 ${1\over
2\pi}\int_C\Tr\,\phi\wedge (-d_A\epsilon)= \Tr\,\gamma \epsilon(p)$，这里我们用了分部积分以及 $d_A\phi=-2\pi \gamma \delta_p$。所以当 $\gamma\not=0$ 时 $\omega_K$ 的上同调类非零。类似地，当 $\beta\not=0$ 时 $\omega_J$ 的上同调类非零。我们将在第~\ref{sec:3.6}~节描述所有这些上同调类。

\bigskip\noindent{\it $\mathbb{L}$ 型的推广}

我们可以容易地把这一点推广到包含一般列维子群 $\mathbb{L}$ 的情形。

我们把参数 $(\alpha,\beta,\gamma)$ 限制为 $\mathbb{L}$-不变的。把 $\EUG_p$ 定义为在 $p$ 处的限制落在 $\mathbb{L}$ 中（而非如上文那样落在 $\TT$ 中）的规范变换群。完全照上面进行，超 Kähler 商 $\EUW(\alpha,\beta,\gamma;p)/\neg/\neg/\EUG_p$
现在给出我们将称为 $\M_{H,\mathbb{L}}$ 的对象，即具有 $\mathbb{L}$ 型分歧结构的希格斯丛的模空间。严格的构造见 \cite{ref23}。当我们不明确指出分歧型时，就意味着取 $\mathbb{L}=\TT$。

在定义 $\M_{H,\mathbb{L}}$ 时必须小心，要求解偏离由 $(\alpha,\beta,\gamma)$ 决定的渐近形式的项比 $1/r\ln r$ 更小，而不仅仅是比 $1/r$ 更小。这样做的目的是排除下文式~\eqref{eq:3.101} 所示的那种渐近行为。

我们之所以为 $\M_{H,\mathbb{L}}(\alpha,\beta,\gamma;p)$ 引入特殊记号，是因为（希望在第~\ref{sec:3.6}~节会变得清楚）它并不是 $\M_H(\alpha,\beta,\gamma;p)$（对正则三元组 $(\alpha,\beta,\gamma)$ 定义）在 $(\alpha,\beta,\gamma)$ 趋于 $\mathbb{L}$-不变值时的极限。这两个空间连维数都不相同。它们的关系是：当三元组 $(\alpha,\beta,\gamma)$ 变为 $\mathbb{L}$-不变时，$\M_{H,\mathbb{L}}$ 是 $\MH$ 的一个奇点轨迹。

本定义的一个特例是 $\mathbb{L}=G$。$G$ 型的分歧结构可以根本没有分歧，因为在 $p$ 处取值于 $G$ 的规范变换群就是全体规范变换。不过，正如第~\ref{sec:2.6}~节末尾所讨论的，若 $G$ 有非平凡的中心 ${\mathcal Z}$，我们的定义有一个轻微的推广。确实，三元组 $(\alpha,\beta,\gamma)
\in \TT\times \mathfrak t\times \mathfrak t$ 是 $G$-不变的，当且仅当 $\beta=\gamma=0$ 且 $\alpha$ 是 ${\mathcal Z}$ 的一个元素。于是 $\M_{H,G}$ 是由 ${\mathcal Z}$ 标记的若干分支之并。其中一个分支（对应 $\mathcal Z$ 的单位元）就是通常的、在 $p$ 处无奇点的希格斯丛的模空间 $\MH$。

另一个极端是 $\mathbb{L}=\TT$。此时 $\M_{H,\mathbb{L}}(\alpha,\beta,\gamma;p)$ 就是我们通常简记为 $\M_H(\alpha,\beta,\gamma;p)$ 的对象。当我们不指定 $\mathbb{L}$ 而谈论带分歧的希格斯丛时，所指的就是这个一般情形。

我们的观点是：支撑在曲面 $D$ 上的面算符（surface operator）由列维型群 $\mathbb{L}$ 的选择来定义。一个特定的面算符导致一个特定的量子场论问题，以及一个特定的带分歧希格斯丛的模空间 $\M_{H,\mathbb{L}}$。在 $\M_{H,\mathbb{L}}$ 的理论中，有时会强调抛物结构（parabolic structure），这一概念我们将在第~\ref{sec:3.4}~节解释。正如我们将看到的（也可以从第~\ref{sec:2.6}~节抛物子群的定义预计到），用抛物结构所作的诠释依赖于 $\alpha$ 取 $\mathbb{L}$-正则值。文献 \cite{ref24} 指出，在保持 $\mathbb{L}$ 固定不变的情况下，通过同时改变 $\beta$、$\gamma$ 和 $\alpha$，可以不遇到奇点而从一种抛物型过渡到另一种抛物型。因此同一个面算符可以导致不同种类的抛物结构；与其按抛物型，不如按列维型来标记面算符。这就是为什么我们宁愿说带分歧的希格斯丛的模空间，而不说抛物希格斯丛（parabolic Higgs bundle）的模空间。不过，当强调复结构 $I$ 的作用时，我们仍会说抛物希格斯丛，因为在那里这一术语是标准的。

\subsection{复观点}
 %% ==== 原文 §3.2：PDF 第 36 页（印刷第 35 页）====
\label{sec:3.2}

人们常常可以通过聚焦于其中一个复结构（例如 $I$）来更多地了解超 Kähler 商 \cite{ref49}。函数
 $\nu_I=\mu_J+i\mu_K$ 在复结构 $I$ 下全纯（holomorphic）。超 Kähler 商 $\EUW/\neg/\neg/\EUG_p$ 等同于用 $\EUG_p$ 对 $\nu_I^{-1}(0)$ 作辛商（symplectic quotient）。（该辛商的定义是：把限制在 $\nu_I^{-1}(0)$ 上的通常矩映射
$\mu_I$ 置零，然后除以 $\EUG_p$。此时 $\nu_I$ 与 $\mu_I$ 都已置零，正给出超 Kähler 商。）按照辛商与复解析商（complex analytic quotient）之间的通常关系，用 $\EUG_p$ 对 $\nu_I^{-1}(0)$ 作辛商，也可以作为几何不变量理论商 $\nu_I^{-1}(0)/\EUG_{p,\C}$ 得到，其中 $\EUG_{p,\C}$ 是 $\EUG_p$ 的复化（complexification）。（由于
$\EUW$ 是带有 $\EUG_p$ 线性作用的仿射空间，一旦选定 $\EUW$ 的某个复结构，$\EUG_p$ 的复化 $\EUG_{p,\C}$ 就自然地作用在 $\EUW$ 上。）若把复结构 $I$ 换成构成 $\EUW$ 超 Kähler 结构的任一复结构，所有这些陈述都有类比。

在目前的情形下，$\nu_I^{-1}(0)/\EUG_{p,\C}$ 可以有全纯的描述，即作为稳定抛物希格斯丛 $(E,\varphi)$ 的模空间。这一点在 \cite{ref22,ref23} 中有解释。就我们的目的而言，并不真正需要这些细节，但确实需要这样一个描述存在这一事实。粗略地说，在指定点 $p\in C$ 处带抛物结构的抛物希格斯丛，是一个希格斯丛
$(E,\varphi)$，其中 $E$ 是 $C$ 上的全纯 $G_\C$-丛，$\varphi$
是 $K_C\otimes {\rm ad}(E)$ 的一个截面，在 $p$ 之外全纯，并在 $p$ 处有一个满足某条件的单极点（simple pole）。若 $\sigma=\half(\beta+i\gamma)$ 是正则的，即它恰好与 $\TT$（更一般地，与 $\mathbb{L}$）交换，则该条件可从式~\eqref{eq:2.11}\ 预计到：$\varphi$ 的极点部分共轭于 $\sigma\,dz/z$。（这里 $z$ 是 $p$ 附近的局部全纯参数。）若 $\sigma$ 不正则，会发生什么则推迟到第~\ref{sec:3.3}~节解释。

抛物希格斯丛的稳定（stable）条件一般依赖于 $\alpha$，但若 $\sigma$ 是一般的，则与 $\alpha$ 无关 \cite{ref24}。在 Kähler 几何中，人们会预期在 $\alpha$ 的特殊值处出现跨墙现象，但在超 Kähler 几何中，对一般的 $\sigma$ 这一现象被回避，正如我们在第~\ref{sec:3.6}~节所解释的。

稳定抛物希格斯丛的模空间（我们暂时记作 $\tilde\M_H$）是复辛流形，带有全纯辛形式（holomorphic symplectic form）$\Omega_I=\omega_J+i\omega_K$。此外，$\tilde\M_H$ 的复结构以及 $\Omega_I$ 的上同调类明显地关于 $\beta+i\gamma$ 全纯。若 $\sigma$ 是一般的，它们也明显地与 $\alpha$ 无关（因为 $\alpha$ 只影响稳定性条件）。

由于把 $\MH(\alpha,\beta,\gamma;p)$ 视为复结构 $I$ 下的复流形时与 $\tilde\M_H$ 相同，所有这些陈述对 $\MH(\alpha,\beta,\gamma;p)$ 都有直接的推论。它们解释了第~\ref{sec:2.2}~节表格中所作的某些论断。特别地，$\MH$ 上的复结构 $I$ 关于 $\beta+i\gamma$ 全纯地变化，且与 $\alpha$ 无关。此外，若以 $[\,\omega]$ 记闭微分形式 $\omega$ 的上同调类，则上同调类 $[\Omega_I]$ 关于 $\beta+i\gamma$ 全纯地变化，且与 $\alpha$ 无关。特别地，类 $[\,\omega_J]$ 与 $[\,\omega_K]$ 与 $\alpha$ 无关。

\bigskip\noindent{\it 其他复结构中的类比}

当然，在复结构 $J$ 中有类似的故事。商 $\nu_J^{-1}(0)/\EUG_{p,\C}$ 可以有全纯的解释，即作为稳定抛物（或带滤过的）局部系统（parabolic (or filtered) local system）的模空间。抛物局部系统是 $C\backslash p$（即去掉点 $p$ 的 $C$）上取值于 $G_{\C}$ 的平坦联络，并带有对绕 $p$ 的单值性的约束。我们把这个单值性记作 $V$。若 $G_\C$ 中的元素 $U=\exp(-2\pi(\alpha-i\gamma))$ 是正则的（意思是与它交换的 $G_\C$ 的子群恰好是环面 $\TT$），则对单值性的约束是 $V$ 必须共轭于 $U$，正如我们从式~\eqref{eq:2.12} 所预期的。否则就需要更精细的描述，我们将它推迟到第~\ref{sec:3.3}~节。抛物局部系统的稳定条件一般依赖于 $\beta$，但若 $\gamma+i\alpha$ 是一般的，则与 $\beta$ 无关 \cite{ref24}。

稳定抛物局部系统的模空间 $\hat \M_H$ 是复辛流形，带有全纯辛形式
$\Omega_J=\omega_K+i\omega_I$。现在我们作与复结构 $I$ 中相同的论证。$\MH$ 视为复结构 $J$ 下的复流形时与 $\hat \M_H$ 相同这一事实，使第~\ref{sec:2.2}~节表格中所作的某些论断变得显然。特别地，$\MH$ 上的复结构 $J$ 关于 $\gamma+i\alpha$ 全纯地变化，且与 $\beta$ 无关。类似地，上同调类 $[\Omega_J]$ 关于 $\gamma+i\alpha$ 全纯，且与 $\beta$ 无关。因此 $[\,\omega_K]$ 与 $[\,\omega_I]$ 与 $\beta$ 无关。


当然，我们可以在复结构 $K$ 中作同样的推理，以证明 $\MH$ 的复结构 $K$ 与辛结构 $\Omega_K$ 关于 $\alpha+i\beta$ 全纯且与 $\gamma$ 无关。于是特别地，上同调类 $[\,\omega_I]$ 与 $[\,\omega_J]$ 与 $\gamma$ 无关。至此完成了对第~\ref{sec:2.2}~节表格中诸论断的论证。




\bigskip\noindent{\it 上同调类的线性性}

利用同样的推理适用于构成 $\MH$ 超 Kähler 结构的任一复结构这一事实，我们可以走得更远。这样的复结构具有形式 $\hat I=pI+qJ+rK$，其中实参数 $p,q,$ 和 $r$ 满足
$p^2+q^2+r^2=1$。第~\ref{sec:2.2}~节表格中的所有论断对任一 $\hat I$ 都有类比。确实，若我们对由复结构 $I,J,K$ 张成的空间作一个 $SO(3)$ 旋转，同时对由三个辛结构 $\omega_I$、$\omega_J$ 和 $\omega_K$ 张成的空间以及由三个变量 $\alpha,$ $\beta$ 和 $\gamma$ 张成的空间作同样的旋转，这些论断仍然成立。这个 $SO(3)$ 不是 $\MH$ 的对称性，但它是我们推导表格中论断所用推理的对称性。

这使我们能对上同调类 $[\,\omega_I]$、$[\,\omega_J]$ 和 $[\,\omega_K]$ 作出一些简单的推断。把这些上同调类记作 $x_i$，$i=1,2,3$。于是 $x_i$ 取值于向量空间
$H^2(\MH,\R)$。它们是 $\alpha,\beta$ 和 $\gamma$ 的函数，我们把这些变量记作 $y_i$，$i=1,2,3$。那么迄今为止所确立的是
% 式 (3.5)
\setcounter{equation}{4}
\begin{equation}\label{eq:3.5}
  {\partial x_i\over\partial
y_j}=0,~~i\not= j.
\end{equation}
然而，通过在一般的复结构 $\hat I$ 中作同样的论证，我们得知：对 $\vec
x=(x_1,x_2,x_3)$ 与 $\vec y=(y_1,y_2,y_3)$ 作一个 $SO(3)$ 旋转后，\eqref{eq:3.5}\ 仍然成立。这意味着
% 式 (3.6)
\setcounter{equation}{5}
\begin{equation}\label{eq:3.6}
  {\partial x_i\over\partial y_j}=\delta_{ij}v
\end{equation}
对
某函数 $v$ 成立（它取值于 $H^2(\MH,\R)$）。\footnote{
这等价于如下陈述：对任意复结构 $\hat I=pI+qJ+rK$，相应的全纯 $(2,0)$ 形式 $\Omega_{\hat I}$（由
$\Omega_I$ 经适当的
$SO(3)$ 旋转得到）的上同调类与 Kähler 形式（它是
$p\alpha+q\beta+r\gamma$ 的倍数）无关。} 再求导一次，由式~\eqref{eq:3.6} 得到
% 式 (3.7)
\setcounter{equation}{6}
\begin{equation}\label{eq:3.7}
  \delta_{ij}{\partial
v\over \partial x_k}=\delta_{ik}{\partial v\over\partial x_j}.
\end{equation}

但这蕴含着（考虑 $i=j$、$i\not= k$ 的情形）对一切 $k$ 有 $\partial v/\partial x_k=0$，故
$v$ 是常数。

\begin{yizhu}[含糊点拨]
由式~\eqref{eq:3.6} 得到~\eqref{eq:3.7} 这一步，原文只说“再求导一次”，其实隐含了链式法则：对 (3.6) 关于 $y_k$ 求导给出 $\partial^2x_i/\partial y_j\partial y_k=\delta_{ij}\partial v/\partial y_k$，由混合偏导的对称性得 $\delta_{ij}\partial v/\partial y_k=\delta_{ik}\partial v/\partial y_j$；再借 $\partial x_i/\partial y_j=\delta_{ij}v$（即 $\partial/\partial y_k=v\,\partial/\partial x_k$，要求 $v\not=0$）把 $\partial v/\partial y$ 换成 $\partial v/\partial x$，方得 (3.7)。因此 $x$ 与 $y$ 之间变换的可逆性与 $v\not=0$ 是隐含前提；若某个参数值上 $v=0$，需用连续性延拓过去。
\end{yizhu}

既然 $v$ 是常数，\eqref{eq:3.6} 就表明 $x_i$ 是 $y_j$ 的线性函数。特别地，$[\,\omega_I]$——我们已知它只依赖于 $\alpha$——实际上是 $\alpha$ 的{\it 线性}函数。于是
% 式 (3.8)
\setcounter{equation}{7}
\begin{equation}\label{eq:3.8}
  [\,\omega_I]=a+\Tr\,\alpha h
\end{equation}
其中 $a$、$h$ 为某些常数。（$a$ 取值于
$H^2(\MH,\R)$，$h$ 取值于 $H^2(\MH,\R)\otimes {\mathfrak t}$。）类似地，我们已经知道 $[\,\omega_J]$ 只依赖于 $\beta$
且在 $\beta=0$ 时为零；现在可由 \eqref{eq:3.6} 推出
% 式 (3.9)
\setcounter{equation}{8}
\begin{equation}\label{eq:3.9}
  [\,\omega_J]=\Tr\,\beta h
\end{equation}
其中 $h$ 与 \eqref{eq:3.8} 中的相同。同理，
% 式 (3.10)
\setcounter{equation}{9}
\begin{equation}\label{eq:3.10}
  [\,\omega_K]=\Tr\,\gamma h.
\end{equation}
我们将在第~\ref{sec:3.5}~节和第~\ref{sec:3.6}~节描述 $a$ 与 $h$。

\bigskip\noindent{\it $\MH$ 作为辛簇}

  正如我们刚才所见，上同调类 $[\,\omega_K]$ 与 $\alpha$ 和 $\beta$ 无关。若 $\MH$ 是紧的，由此便可得：以 $\omega_K$ 为辛结构的实辛簇（real symplectic variety）$\MH$ 与 $\alpha$、$\beta$ 无关。由于 $\MH$ 非紧，这一结论不能仅由辛形式上同调类的常性推出；但更仔细地考察我们用来证明该常性的事实，就可以得到它。

为了在改变 $\alpha$ 时不改变实辛簇 $(\MH,\omega_K)$，我们只需把 $\MH$ 放在复结构 $I$ 中看。由于有抛物希格斯丛的解释，这个复结构以及相应的全纯二形式
$\Omega_I=\omega_J+i\omega_K$ 与 $\alpha$ 无关，只要避开某些奇点。于是我们可以在保持实辛簇
$(\MH,\omega_K)$ 不变的情况下改变 $\alpha$。

类似地，为了在改变 $\beta$ 时不改变实辛簇 $(\MH,\omega_K)$，我们只需把 $\MH$ 放在复结构
$J$ 中看。由于有带滤过局部系统的解释，这个复结构以及相应的全纯二形式
$\Omega_J=\omega_K+i\omega_I$ 与 $\beta$ 无关，只要避开某些奇点。于是我们可以在保持实辛簇
$(\MH,\omega_K)$ 不变的情况下改变 $\beta$。

此过程中必须避开的奇点在第~\ref{sec:3.6}~节中描述，其实余维（codimension）至少为二。于是我们断言：实辛簇 $(\MH,\omega_K)$ 与 $\alpha$ 和 $\beta$ 无关——这是应用于几何朗兰兹纲领的一个重要结果。

一个要点是：该论证给出的同构不是典范的，而依赖于改变 $\alpha$ 与 $\beta$ 的路径。原因很简单：改变 $\alpha$ 时我们保持一个结构不变，而改变 $\beta$ 时保持的是另一个结构不变。若同时改变 $\alpha$ 和 $\beta$，则除 $\omega_K$ 之外没有结构被固定，$\MH$ 的变化由一个辛同胚（symplectomorphism）实现。
让 $\alpha$ 和 $\beta$ 沿避开奇点的闭合回路变化，一般会得到 $\MH$ 的一个“拓扑上”非平凡的辛同胚，即它不能通过一族辛同胚形变到恒等映射。我们将在第~\ref{sec:3.7}~节和第~\ref{sec:4.5}~节系统地展开这一思想，以得到仿射辫群（affine braid group）在 $\MH$ 上的膜（brane）上的作用。

\bigskip\noindent{\it $\C^*$ 作用}

旋转 $\MH$ 的复结构空间或辛结构空间的那个 $SO(3)$ 群不是 $\MH$ 的对称性。然而它的一个子群是对称性，正如对无抛物结构的通常希格斯丛那样，后者在 \cite{ref21} 第 107--8 页中有描述。

先考虑 $\beta=\gamma=0$ 的情形。让群
${\mathcal U}_1\cong U(1)$ 作用在 $(A,\phi)$ 上：保持 $A$ 不变，而把 $\varphi\to\lambda\varphi$，其中 $|\lambda|=1$。由于 $\phi=\varphi+\bar\varphi$（其中 $\varphi$ 与 $\bar\varphi$
分别为 $(1,0)$ 型和 $(0,1)$ 型），这决定了 $\phi$ 的变换：
% 式 (3.11)
\setcounter{equation}{10}
\begin{equation}\label{eq:3.11}
  \phi\to\lambda \varphi+\bar\lambda
\bar\varphi.
\end{equation}

   ${\mathcal U}_1$ 的作用保持奇点的刻画式~\eqref{eq:3.13} 不变（只要 $\beta=\gamma=0$）。它保持矩映射 $\mu_I$ 不变，同时旋转
$\mu_J-\mu_K$ 平面。因此它给出超 Kähler 矩映射构造的一个明显对称性，以及 $\MH$ 上超 Kähler 度规的一个对称群。事实上，${\mathcal U}_1$ 是旋转三个复结构的那个 $SO(3)$ 的 $SO(2)$ 子群。它作用在复结构族 $\hat I=pI+qJ+rK$ 上：保持 $p$ 不变，旋转
$q-r$ 平面。或者，若用复变量 $w$ 参数化复结构族，令

% 式 (3.12)
\setcounter{equation}{11}
\begin{equation}\label{eq:3.12}
  I_w={1-\bar w w\over 1+\bar w
w}I+{i(w-\bar w)\over 1+\bar w w}J+{w+\bar w\over 1+\bar w w}K,
\end{equation}

则 $\C^*$ 通过 $w\to\lambda^{-1}w$ 作用在参数 $w$ 上。

\begin{yizhu}[易错点]
此处引用的式~\eqref{eq:3.13} 在第~\ref{sec:3.3}~节（PDF 第 41 页）才出现，属于前向编号引用；它与式~\eqref{eq:3.2} 是同一奇点条件的重述。原文中这类前向引用并不罕见（例如本章第~\ref{sec:3.1}~节末尾也前向引用了式~\eqref{eq:3.101}），故不宜视为笔误；译文照录原编号，读者若在本节找不到 (3.13)，可直接按式~\eqref{eq:3.2} 理解。
\end{yizhu}

若放宽条件 $|\lambda|=1$，我们便不再得到 $\MH$ 超 Kähler 度规的对称性。然而，完全如 \cite{ref21} 中那样，我们得到一个群 ${\mathcal U}\cong \C^*$，它作用在 $\MH$ 上，保持复结构 $I$ 不变，并使复结构族 $I_w$ 按 $w\to\lambda^{-1}w$ 变换，与 $|\lambda|=1$ 时相同。不动点为 $w=0$（$I_w=I$）和 $w=\infty$（$I_w=-I$）。其余所有复结构
$I_w$ 在 ${\mathcal U}$ 的作用下彼此等价。

这一切都可推广到 $\beta,\gamma\not=0$ 的情形，只需把参数 $\beta$ 和 $\gamma$ 按 $(\beta+i\gamma)\to\lambda(\beta
+i\gamma)$ 变换。对 $|\lambda|=1$，这一点可由如下观察得到：若 $\beta$ 和 $\gamma$ 按所述方式变换，则变换~\eqref{eq:3.11} 保持奇点~\eqref{eq:3.13} 不变。

\subsection{非正则情形}
 %% ==== 原文 §3.3：PDF 第 41 页（印刷第 40 页）====
\label{sec:3.3}

现在我们来描述：当数对 $(\alpha,\gamma)$ 或数对 $(\beta,\gamma)$ 非正则时，上述某些论断会发生什么变化。这或许有助于读者理解 \cite{ref22} 中的构造，并且是几何朗兰兹纲领某些应用的有用背景。

\bigskip\noindent{\it 复结构 $J$}

我们从复结构 $J$ 开始，其中希钦方程的解对应于抛物局部系统。另外，为简单起见，先考虑用来定义面算符的列维子群为 $\mathbb{L}=\TT$ 这一基本情形。

考虑在点 $p\in C$ 处具有孤立奇点的希钦方程之解。按第~\ref{sec:2.1}~节的方式在 $p$ 附近取坐标（于是 $p$ 由 $z=0$ 定义，其中 $z=x_1+ix_2=re^{i\theta}$），并假定解在 $r=0$ 附近的行为为

% 式 (3.13)
\setcounter{equation}{12}
\begin{equation}\label{eq:3.13}
  \begin{aligned}A & = \alpha \,d\theta +\dots\\
           \phi & = \beta\,{dr\over r}-\gamma\,d\theta+\dots,\end{aligned}
\end{equation}

其中省略号表示当 $r\to 0$ 时比 $1/r$ 更不奇异的项。希钦方程保证取值于 $G_\C$ 的联络 $\CA=A+i\phi$ 是平坦的。

对 $b>0$（但足够小，使得坐标 $r,\theta$ 在 $r\leq b$ 上有定义），令 $C_b$ 为圆周 $r=b$，并令 $V_b$
为平坦联络 $\CA$ 绕 $C_b$ 的单值性。$V_b$ 有两条基本性质：(1) $V_b$ 的共轭类（conjugacy class）与 $b$ 无关。这是因为联络是平坦的。既然我们只关心单值性的共轭类，就任取一个 $V_b$ 并称之为 $V$，即绕奇点的单值性。(2) 若令
$U=\exp\left(-2\pi(\alpha-i\gamma)\right)$，则

% 式 (3.14)
\setcounter{equation}{13}
\begin{equation}\label{eq:3.14}
  \lim_{b\to 0}V_b=U.
\end{equation}
这是因为在 $b\to 0$ 的极限下，只需由 $A$ 与 $\phi$ 中最奇异的项就能算出和乐（holonomy），而这样算出的结果就是 $U$。

若 $U$ 是正则的，这两条性质便蕴涵 $V$ 共轭于 $U$。对 $SL(2,\C)$，可如下证明。两条性质蕴涵
$\Tr \,U=\Tr\, V$（迹在二维表示中取）；若 $U$ 是正则的，这就蕴涵 $U$ 与 $V$
共轭。类似的推理对任意 $G$ 都成立，只需把迹换成一整套不变函数。若 $U$ 不正则，$U$ 与 $V$ 未必共轭。例如对
$G=SL(2,\C)$，考虑非正则元素 $U=1$，它对应于 $\alpha=\gamma=0$。对 $U=1$，可以用下式满足上述两个条件
% 式 (3.15)
\setcounter{equation}{14}
\begin{equation}\label{eq:3.15}
  V_b=\left(\begin{matrix}1& b\\ 0 &
1\end{matrix}\right).
\end{equation}
  $b>0$ 时的 $V_b$ 全都彼此共轭，且
$\lim_{b\to 0}V_b=1$。这种形式的元素 $V$ 称为幂幺的（unipotent）（这简单地意味着 $V-1$ 是幂零的（nilpotent））。

因此，$\alpha=\gamma=0$、$\beta\not=0$ 的希钦方程之解可以具有幂幺但不等于 1 的单值性。这不仅可能，而且是一般行为，理由很简单：$V$ 为幂幺只需要一个复条件（可表述为
$\Tr\,V=2$），而 $V$ 等于 1 需要三个复条件。比较这些维数，人们或许会以为：若
$\alpha=\gamma=0$，则那些 $V$ 实际为 1 的希格斯丛将构成复余维为二的一族。更仔细的分析表明，当 $\beta=0$ 时（此时 $V=1$ 的希格斯丛是 $A_1$ 轨形奇点的轨迹）这是对的；但当 $\beta\not=0$ 时，$V=1$ 的轨迹被“爆破”（blown up），其复余维为一。关于一般行为、奇点和爆破的论断在下文、特别是第~\ref{sec:3.6}~节中应当会变得更清楚。爆破的描述方式是：指定 Simpson \cite{ref22} 所谓的局部系统的“滤过”（filtration）。

设 $\mathfrak C$ 是 $G_\C$ 中包含 $U$ 的共轭类。对任意 $G$ 与 $\mathfrak C$ 的任意选择，都存在一有限组共轭类
$\mathfrak C_\lambda$（$\lambda=1,\dots, s$），具有如下性质：一族元素 $V_b\in {\mathfrak C}_\lambda$ 可以在 $b\to
0$ 时以 $\mathfrak C$ 中的元素为极限。换言之，$\mathfrak C$ 位于 $\mathfrak C_\lambda$ 的闭包中。我们称这些 ${\mathfrak C}_\lambda$ 为从属于（affiliated to）$\mathfrak C$ 的共轭类。例如，若 $U$ 正则，则唯一从属的共轭类就是 ${\mathfrak C}$ 本身。在另一极端，若 $U=1$，则从属的共轭类恰是那些参数化 $G$ 的幂幺元素的类。对秩较大的 $G$，这样的类有很多（对 $SL(N,\C)$ 而言，由对角线上全为 1、各块大小不同的分块三角矩阵给出）。

无论如何，总存在唯一的维数最大的从属共轭类 ${\mathfrak C}^*$，其维数恰为
${\rm dim}\,G-r$，其中 $r$ 为 $G$ 的秩（rank）。它称为正则共轭类，因为它参数化的元素是正则的，即只与 $G_\C$ 的一个 $r$ 维子群交换。

类 $\mathfrak C$ 本身未必正则，但它有另一个区别性性质。在从属的共轭类中，$\mathfrak C$ 是唯一“半单”（semi-simple）的，即它参数化的群元素是半单的（可以共轭到极大环面（maximal torus）中）。$\mathfrak C$ 当然是半单的，因为它包含 $U=\exp(-2\pi(\alpha-i \gamma))$，而后者是 $\TT$ 的元素。


但若 $U$ 不正则，则 ${\mathfrak C}\not={\mathfrak C}^*$，而从属于 $U$ 的正则共轭类中的元素不是半单的。例如，若 $G=SU(N)$ 且 $U=1$，则从属于 $U$ 的正则共轭类包含“主幂幺元”（principal unipotent element）

% 式 (3.16)
\setcounter{equation}{15}
\begin{equation}\label{eq:3.16}
  V=\left(\begin{matrix}1& 1& 0&\dots&0\\
                         0&1&1&\dots&0\\
                           & & &\vdots&\\
                           0&0&0&\dots&1\end{matrix}\right),
\end{equation}

其主对角线上及紧邻主对角线上方全为 $1$，其余位置为零。对任意 $\alpha$ 和 $\gamma$，一般的希格斯丛给出的局部系统，其单值性落在与 $U$ 相关联的正则共轭类中。

\bigskip\noindent{\it 复结构 $I$}

现在考虑复结构 $I$，其中希钦方程的解是抛物希格斯丛。

这里先详细考虑 $SL(2,\C)$ 的例子会有帮助。取 $\beta=\gamma=0$（这是 $SL(2,\C)$ 中 $\beta$ 和 $\gamma$ 唯一非正则的取法）；又因我们希望假定 $\alpha$ 正则，故取

% 式 (3.17)
\setcounter{equation}{16}
\begin{equation}\label{eq:3.17}
  \alpha=iy\left(\begin{matrix}1 & 0\\ 0&-1\end{matrix}\right),
\end{equation}

其中 $0<y<1/2$。（因子 $i$ 的理由是：由于 $\alpha$ 取值于 $SU(2)$ 的实李代数，它在酉表示（unitary representation）中是反厄米的（anti-hermitian）。）因此当 $r\to 0$ 时，希钦方程之解的极限形式为

% 式 (3.18)
\setcounter{equation}{17}
\begin{equation}\label{eq:3.18}
  \begin{aligned}A&=iy\,d\theta\left(\begin{matrix}1 & 0\\ 0&-1\end{matrix}\right)\\
                   \phi&=0.\end{aligned}
\end{equation}

我们想知道：当把式~\eqref{eq:3.18} 换成一个具有同样 $r\to 0$ 渐近行为、但更一般的希钦方程之解时，$\varphi$（定义为 $\phi$ 的 $(1,0)$ 部分）能获得何种极点。

乍看之下，似乎根本不可能有极点。按定义，我们要用比 $1/r$ 更不奇异的项来扰动极限解~\eqref{eq:3.18}。于是似乎 $\varphi$ 的奇异性不足以产生极点。然而，要判断 $\varphi$ 是否有极点，需要把丛 $E$ 在奇点 $p$ 附近的全纯结构平凡化（trivialize）。结果表明，一旦这样做，$\varphi$ 就可以有极点。

把希钦方程在解~\eqref{eq:3.18} 附近展开时，线性化方程有如下解

% 式 (3.19)
\setcounter{equation}{18}
\begin{equation}\label{eq:3.19}
  \phi=\epsilon {dz\over z}(\bar z z)^{y}\left(\begin{matrix}0&1\\ 0&0\end{matrix}\right),
\end{equation}

其中 $\epsilon$ 是小参数。该解比 $1/r$ 更不奇异，因此包含这一扰动与渐近行为~\eqref{eq:3.18} 相容。（当然，为了得到关于 $\phi$ 的实解，必须减去厄米共轭的解。$\phi$ 的 $(1,0)$ 部分必须是上三角的（upper triangular），否则一个类似的解在 $r=0$ 处会比 $1/r$ 更奇异。）

为把 $r=0$ 附近的全纯丛 $E$ 平凡化，写出定义 $E$ 的全纯结构的适当的
 $\bar\partial_A$ 算符：

% 式 (3.20)
\setcounter{equation}{19}
\begin{equation}\label{eq:3.20}
  \bar\partial_A=d\bar z\left({\partial\over\partial\bar
z}+A_{\bar z}\right)= d\bar z\left({\partial\over\partial\bar
z}-{y\over 2\bar z}\left(\begin{matrix}1&0\\ 0&-1\end{matrix}\right) \right).
\end{equation}
我们可以把它写成
% 式 (3.21)
\setcounter{equation}{20}
\begin{equation}\label{eq:3.21}
  \bar\partial_A=f\bar\partial f^{-1},
\end{equation}

其中 $\bar\partial =d\bar z\partial/\partial\bar z$ 是标准的
$\bar\partial$ 算符，而
% 式 (3.22)
\setcounter{equation}{21}
\begin{equation}\label{eq:3.22}
  f=\left(\begin{matrix}(\bar z
z)^{y/2}& 0\\ 0& (\bar z z)^{-y/2}\end{matrix}\right).
\end{equation}
于是若 $\bar\partial_A\varphi=0$（这是希钦方程的一部分），则 $\bar\partial(f^{-1}\varphi f)=0$。但
% 式 (3.23)
\setcounter{equation}{22}
\begin{equation}\label{eq:3.23}
  f^{-1}\varphi
f= \epsilon {dz\over z}\left(\begin{matrix}0&1\\ 0&0\end{matrix}\right).
\end{equation}
所以结论是：相对于丛 $E$ 的平凡化，$\varphi$ 可以有极点，但该极点的留数（residue）是严格上三角的（strictly upper triangular）（在使 $-i\alpha$ 对角、且本征值沿对角线递减的基下）。

在上述讨论中，我们从 $\beta=\gamma=0$ 出发，故 $\sigma={1\over 2}(\beta+i\gamma)$
也为零。结果是 $\varphi$ 中极点的留数为

% 式 (3.24)
\setcounter{equation}{23}
\begin{equation}\label{eq:3.24}
  \tau=\epsilon\left(\begin{matrix}0 & 1\\ 0 & 0 \end{matrix}\right).
\end{equation}

$\tau$ 的共轭类与 $\epsilon$ 无关，且 $\tau$ 在
$\epsilon\to 0$ 时的极限为 $\sigma$。这些性质蕴涵：所有不变多项式在 $\sigma$ 与 $\tau$ 处取相同的值；对 $SL(2,\C)$，这简单地意味着 $\Tr\,\sigma^2=
\Tr\,\tau^2=0$。

一般地，对任意 $G$，设 $\mathfrak c$ 是李代数 $\mathfrak g_\C$ 中包含 $\sigma={1\over
2}(\beta+i\gamma)$ 的轨道或共轭类。由于 $\sigma$ 本身是半单的，这个共轭类参数化 $\mathfrak g_\C$ 的半单元素（即可共轭到极大环面的元素）。若 $\mathfrak c_\lambda$ 中的一列元素可以收敛到 $\mathfrak c$ 中的元素，换言之，若 $\mathfrak c$
位于 $\mathfrak c_\lambda$ 的闭包中，则称共轭类 ${\mathfrak c}_\lambda$ 从属于 $\mathfrak c$。一般地，从属于 $\mathfrak c$ 的共轭类只有有限多个。希格斯场极点的留数 $\tau$ 总取值于某个从属的共轭类。这一点可由与上例中所解释的类似的推理得到。


若 $\sigma$ 既是正则的又是半单的，则 $\mathfrak c$ 本身是唯一从属的共轭类，特别地 $\tau$ 共轭于 $\sigma$。在另一极端，若 $\sigma=0$，则从属于 $\mathfrak c$ 的共轭类恰是 $\mathfrak g_\C$ 的幂零元素的类。

对每个 $\mathfrak c$，都存在唯一的维数最大的从属共轭类
${\mathfrak c}^*$，其维数恰为 ${\rm
dim}(G)-r$。它参数化 $\mathfrak g_\C$ 的正则元素，即只与 $G_\C$ 的一个 $r$ 维子群交换的元素。例如，若 $\sigma=0$ 且 $G=SU(N)$，则从属的正则共轭类包含元素 $v=V-1$，其中 $V$ 由式~\eqref{eq:3.16} 定义。对任意 $\sigma$，一般的抛物希格斯丛的极点留数 $\tau$ 都落在正则从属共轭类 ${\mathfrak c}^*$ 中。

上述内容大部分可以概括为：李代数
$\mathfrak g_\C$ 中的共轭类在许多相关方面表现得与群 $G_\C$ 中的共轭类相似。

\bigskip\noindent{\it 重新表述}


回到 $SU(2)$ 的例子：我们关于 $\beta=\gamma=0$ 时 $\varphi$ 极点行为的结果可表述如下。对 $G=SU(2)$，唯一可能的列维子群是 $\mathbb{L}=\TT$，这正是上面讨论的情形。$\alpha$ 的选择决定一个抛物子群 $\EUP$ 与一个幂幺根（unipotent radical）$\EUN$。在上例中，$\EUP$ 是上三角矩阵构成的群，而 $\EUN$ 是这样一个群：其李代数 $\mathfrak n$ 由严格上三角矩阵构成。结果~\eqref{eq:3.23} 表明 $\varphi$ 的极点留数取值于 $\mathfrak n$。造成这一点的机制很简单：为了使 $\phi$（在微分几何描述中）比 $1/r$ 更不奇异，它在全纯描述中的极点留数 $\tau$ 必须满足 $-i[\alpha,\tau]=\lambda\tau$，$\lambda>0$。对更高秩的 $G$（仍假定 $\beta=\gamma=0$），相应的类比是：$\tau$ 必须是 $\mathfrak g$ 中满足此条件的元素的线性组合。换言之，$\tau$ 取值于 $\mathfrak n$。这对任意列维子群 $\mathbb{L}$ 的选择与任意 $\mathbb{L}$-正则的 $\alpha$ 都成立。

回到 $SU(2)$，在不改变 $\alpha$ 的情况下，把 $\beta$ 和 $\gamma$ 扰动为非零，例如取 $\half(\beta+i\gamma)={\rm diag}(q,-q)$（其中 $q\in \C$）。此时会怎样？我们必须令 $\phi=\beta (dr/r)-\gamma\,d\theta$
再加上在 $r=0$ 处更不奇异的项。一种可能是 $\phi=\beta(dr/r)-\gamma\,d\theta$ 精确成立。这对应于 $\varphi=(dz/z){\rm diag}(q,-q)$。但正如式~\eqref{eq:3.19} 那样，我们可以对 $\varphi$ 作上三角形变\footnote{对 $q\not=0$，解的细节更为复杂，还必须修改 $A$。适当的解由纳姆方程描述；见第~\ref{sec:3.8}~节。与 $q=0$ 时一样，对 $\varphi$ 作下三角修改会比 $1/r$ 更奇异。}，于是 $\varphi$ 的极点部分的一般可能是
% 式 (3.25)
\setcounter{equation}{24}
\begin{equation}\label{eq:3.25}
  \varphi\sim{dz\over z}\left(\begin{matrix}q&
*\\ 0 & -q\end{matrix}\right).
\end{equation}

记作 $*$ 的右上元素在 $q\not=0$ 时不影响极点留数的共轭类。注意：若 $\varphi$ 与 $\tilde\varphi$ 是带有这类极点的两个希格斯场，则 $\varphi-\tilde\varphi$ 的极点具有 $\mathfrak n$-值留数：

% 式 (3.26)
\setcounter{equation}{25}
\begin{equation}\label{eq:3.26}
  \varphi-\tilde\varphi\sim {dz\over
z}\left(\begin{matrix}0&*\\ 0&0\end{matrix}\right).
\end{equation}


由类似的推理，这不仅对 $SL(2,\C)$ 成立，而且对任意规范群 $G$、任意列维子群 $\mathbb{L}$ 的选择以及任意 $\mathbb{L}$-正则的 $\alpha$ 都成立。一般性的陈述是：$\varphi$ 的极点留数取值于 $\mathfrak p$（即 $\EUP$ 的李代数），并且等于 $\sigma=\half(\beta+i\gamma)$
模去 $\mathfrak n$ 中的一个元素。最后这一陈述可以非正式地概括为：极点留数的“本征值”与 $\sigma$ 的本征值相同。

\bigskip\noindent{\it 再论复结构 $J$}

我们在复结构 $J$ 和 $I$ 中对非正则行为的讨论是不对称的。在复结构 $J$ 中，我们关于共轭类的闭包作了一些较为抽象的推理；而在复
结构 $I$ 中，我们分析了希钦方程扰动的行为。当然，通过分析希钦方程，我们可以更明确地说明复结构 $J$ 中发生的事。这还将使我们获得更多信息。

我们对任意规范群 $G$ 与列维子群 $\mathbb{L}$ 进行讨论。首先取 $\alpha=\gamma=0$，但取 $\beta$ 为一般的，换言之 $\mathbb{L}$-正则的。这对应于希钦方程的一个解
% 式 (3.27)
\setcounter{equation}{26}
\begin{equation}\label{eq:3.27}
  A=0, ~\phi=\beta
{dr\over r}.
\end{equation}

我们把它扰动为

% 式 (3.28)
\setcounter{equation}{27}
\begin{equation}\label{eq:3.28}
  \begin{aligned}A&=a(r)\,d\theta,\\ \phi&=\beta{dr\over r}+c(r)\,d\theta,\end{aligned}
\end{equation}

其中 $a$ 和 $c$ 必须在 $r=0$ 处为零（因为与极限解~\eqref{eq:3.27} 的偏离必须比 $1/r$ 更不奇异），而我们将只保留 $a$ 和 $c$ 的一阶项。所得方程可写成

% 式 (3.29)
\setcounter{equation}{28}
\begin{equation}\label{eq:3.29}
  r{d\over dr}(a+ic)=[-i\beta,a+ic].
\end{equation}

要使 $a+ic$ 在 $r=0$ 处为零，它必须取值于 $\mathfrak n$，即 $\mathfrak g$ 中由满足 $-i[\beta,\psi]=\lambda\psi$（$\lambda>0$）的向量 $\psi$ 张成的子空间。因此，平坦联络 ${\mathcal A}=A+i\phi$ 的单值性（在此近似下为 $U=\exp(-2\pi(a+ic))$）取值于 $\EUN$，即由 $\beta$ 决定的抛物子群的幂幺根。

\begin{yizhu}[易错点]
式~\eqref{eq:3.29} 就是本题中纳姆方程的线性化形式：令 $t=\ln r$，它写成 $d(a+ic)/dt=[-i\beta,a+ic]$，即矩阵值函数按伴随作用随 $t$ 演化。这里出现 $-i\beta$ 而不是 $\beta$，源于复化联络 $\CA=A+i\phi$ 与坐标 $z=re^{i\theta}$ 的约定；与不同文献的纳姆方程对照时，须同时跟踪 $A$ 与 $\phi$ 的角色以及因子 $i$ 的归属，否则会差一个整体符号。脚注~17 指出，$q\not=0$ 时完整解由纳姆方程描述（第~\ref{sec:3.8}~节）；本节只用到如下性质：$\mathfrak n$ 中的修改使 $a+ic$ 在 $r=0$ 处为零，从而使偏离比 $1/r$ 更不奇异。
\end{yizhu}

$\EUN$ 的每个元素都是 $G_\C$ 的幂幺元素。在半单李群的理论中证明了：$\EUN$ 的一般元素落在 $G_\C$ 中一个称为理查森类（Richardson class）${\mathfrak C}_\mathbb{L}$ 的幂幺共轭类中。从定义看，理查森类似乎依赖于 $\EUP$、从而依赖于 $\beta$；但可以证明，理查森类实际上由 $\mathbb{L}$ 决定。（一般地，不同的 $\mathbb{L}$ 可以导致相同的理查森类。）上一段得到的单值性 $U$ 除了取值于 $\EUN$ 之外没有其他约束，因此一般地它取值于共轭类 ${\mathfrak C}_\mathbb{L}$。

例如，若 $\mathbb{L}=\TT$，则 ${\mathfrak C}_\mathbb{L}$ 就是式~\eqref{eq:3.16} 中对 $SL(N,\C)$ 描述的正则幂幺共轭类。使 $U-1$ 为严格上三角的 $N\times N$ 矩阵 $U$ 一般属于这个共轭类。再举一例：对 $SL(3,\C)$，若
$\mathbb{L}$ 由形如

% 式 (3.30)
\setcounter{equation}{29}
\begin{equation}\label{eq:3.30}
  \left(\begin{matrix}* & * & 0 \\
                        * & * & 0 \\
                        0 & 0 & * \end{matrix}\right)
\end{equation}
的矩阵构成，且 $\alpha$ 使得幂幺根 $\EUN$ 由形如

% 式 (3.31)
\setcounter{equation}{30}
\begin{equation}\label{eq:3.31}
  \left(\begin{matrix}1 & 0 &* \\
                        0 & 1 & * \\
                        0 & 0 & 1 \end{matrix}\right),
\end{equation}
的矩阵构成，则 $\EUN$ 的一般元素在 $\mathbb{L}$ 下共轭于

% 式 (3.32)
\setcounter{equation}{31}
\begin{equation}\label{eq:3.32}
  \left(\begin{matrix}1 & 0 & 0 \\
                        0 & 1 & 1 \\
                        0 & 0 & 1 \end{matrix}\right).
\end{equation}
这是理查森共轭类 ${\mathfrak C}_\mathbb{L}$ 的一个代表。一般地，对 $SL(N,\C)$，每个幂幺共轭类都是理查森类，但对其他群并非如此。

上述讨论可以推广到 $\alpha,\gamma\not=0$ 的情形。若 $U=\exp(-2\pi(\alpha-i\gamma))$ 是 $\mathbb{L}$-正则的，则单值性 $V$ 共轭于 $U$。一般地，它具有形式 $V=UN$，其中 $N$ 取值于幂幺根 $\EUN$。等价地说，单值性落在 $\EUP$ 中某个闭包包含 $U$ 的共轭类里。此外，在 $V$ 的共轭类范围内，可以假定 $U$ 和 $N$ 交换。此时一般地，$N$ 就是 $\EUN$ 中与 $U$ 交换的一般元素。

若 $N$ 具有这种形式，则 $\EUP$ 中元素 $V=UN$ 在 $G_\C$ 中的轨道与 $G_\C/\mathbb{L}_\C$ 同维，后者是 $\TT_\C$ 中一般 $\mathbb{L}$-正则元素在 $G_\C$ 中的轨道。我们把 $\EUP$ 中具有这一性质的元素称为 $\mathbb{L}$-正则的；于是特别地，对任意 $U$，复结构 $J$ 中局部系统的一般单值性 $V=UN$ 都是 $\mathbb{L}$-正则的。

\begin{yizhu}[物理图像补充]
在复结构 $J$ 中，面算符的单值性按 $V=UN$ 分解：半单部分 $U=\exp(-2\pi(\alpha-i\gamma))$ 由 $(\alpha,\gamma)$ 决定；幂幺部分 $N$ 落在由 $\beta$ 决定的幂幺根 $\EUN$ 中，并且在 $V$ 的共轭类范围内可与 $U$ 交换。一般情形下 $N$ 遍历 $\EUN$ 中与 $U$ 交换的一般元素，故单值性落在正则从属类（理查森类 ${\mathfrak C}_{\mathbb{L}}$）中。这正对应第~\ref{sec:2.6}~节与第~\ref{sec:3.1}~节的论断：面算符应按列维型而非按抛物型标记——因为 $\beta$、$\gamma$ 可以在不遇奇点的情况下改变抛物结构，而 $\mathbb{L}$ 与 $\alpha$ 的 $\mathbb{L}$-正则性保持不变。换言之，在 $\MH$ 的模空间层面，$\beta$、$\gamma$ 与 $N$ 都是可以连续变化的参数，真正刻画面算符类型的是 $\mathbb{L}$。
\end{yizhu}

例如，理查森轨道都是 $\mathbb{L}$-正则的。在上述 $SL(3,\C)$ 例子中，$G_\C/\mathbb{L}_\C$ 的维数为 4，这也是式~\eqref{eq:3.32} 所述理查森轨道的维数。若给定的幂幺轨道是若干不同列维子群 $\mathbb{L}_i$ 的理查森轨道，则它对每个 $i$ 都是 $\mathbb{L}_i$-正则的。
